\documentclass[11pt]{article}
\usepackage[margin=1in]{geometry}
\usepackage{amsmath,amssymb,amsthm}
\usepackage{booktabs}
\usepackage{listings}
\usepackage{xcolor}
\usepackage[numbers]{natbib}
\usepackage{hyperref}

\lstdefinelanguage{lean}{
  morekeywords={def,theorem,lemma,example,namespace,end,import,where,instance,fun,by,decide,Prop,Type,open,intro,exact,rw,refine,have,simpa,apply,with,induction,using,by_contra,push_neg,obtain,cases,unfold,omega,calc,classical,constructor,linarith,rcases,rintro,subst,ring,simp,only,at},
  sensitive=true, morecomment=[l]{--}, morecomment=[s]{/-}{-/}, morestring=[b]"
}
\lstset{language=lean, basicstyle=\ttfamily\small, breaklines=true, columns=fullflexible,
  keywordstyle=\color{blue!60!black}\bfseries, commentstyle=\color{gray},
  literate={∀}{{$\forall$}}1 {→}{{$\to$}}1 {↔}{{$\leftrightarrow$}}1 {∧}{{$\land$}}1
    {∨}{{$\lor$}}1 {¬}{{$\neg$}}1 {∈}{{$\in$}}1 {≤}{{$\le$}}1 {∏}{{$\prod$}}1
    {ℝ}{{$\mathbb{R}$}}1 {ℕ}{{$\mathbb{N}$}}1 {ι}{{$\iota$}}1 {·}{{$\cdot$}}1 {−}{{$-$}}1
    {⊤}{{$\top$}}1 {⊥}{{$\bot$}}1 {⟨}{{$\langle$}}1 {⟩}{{$\rangle$}}1 {←}{{$\leftarrow$}}1
    {⦃}{{$\{$}}1 {⦄}{{$\}$}}1 {⊢}{{$\vdash$}}1 {∃}{{$\exists$}}1}

\newtheorem{theorem}{Theorem}
\newtheorem{lemma}[theorem]{Lemma}

\title{A route is no stronger than its weakest link,\\ and alternatives never hurt:\\ serial and parallel route reliability with a Lean proof}
\author{Justin Hart\\Viridis LLC}
\date{26 November 2026}

\begin{document}
\maketitle

\begin{abstract}
Earlier notes in this series gave Lean proofs about corridor counts in Hierarchical Dendritic Forest
Management (HDFM). A connected network on $n$ habitat patches needs at least $n-1$ corridors
(\texttt{P6.MinimalNetwork}), and surviving the loss of any single corridor needs at least $n$
(\texttt{P6.SingleFailureRobust}). This note moves from counting corridors to their reliability. A route
is a chain of links crossed with probabilities $p_i\in[0,1]$, and its serial success probability is
$\prod_i p_i$. Alternative routes with success probabilities $q_i\in[0,1]$ combine by the parallel
formula $1-\prod_i(1-q_i)$. We give a Lean~4 proof of the sealed statement \texttt{P6.SerialParallel}.
(1)~The serial product is at most every $p_j$: a route is no stronger than its weakest link.
(2)~The parallel formula lies between every $q_j$ and $1$: adding alternatives never lowers it.
(3)~For two routes, $1-(1-p)(1-q)=\max(p,q)$ exactly when $\min(p,q)=0$ or $\max(p,q)=1$, so a second
route gives a strict gain unless it is dead or the first route is certain. Part~1 is the same inequality
as part~3 of \texttt{P3.CascadedRecovery}, which a run on 16 November 2026 proved. The inequalities are
pure algebra. Reading the parallel formula as the probability that at least one route succeeds requires
the routes to fail independently. That is a field assumption, not part of the mathematics, and it fails
when routes share links or a common disturbance. These are elementary facts, not new theorems. We also
ran the registered test plan \texttt{P6-NucleationIslandSize}. It returned \textsc{not\_estimable}
because its frozen code looks for a column the data file does not contain.
\end{abstract}

\section{Introduction}
Graph models of landscape connectivity represent habitat patches as vertices and functional links
between them as edges~\citep{urbankeitt2001}. The two previous notes in this project (P6-T002, P6-T003)
studied how many corridors a network needs. Counting treats a corridor as either present or absent. In
practice a corridor is crossed with some probability less than one, and a planner wants to know how
those probabilities combine along a route and across alternative routes.

\paragraph{Questions.} (a) How does the reliability of a route of several links compare with its
individual links? (b) How much can alternative routes add, and when does a second route add nothing?

\paragraph{Answers.} (a) The route's serial success probability is at most that of its weakest link.
(b) The combined probability that at least one route succeeds, computed with the independence formula,
is at least that of the best single route and at most $1$. For two routes it equals the better route
exactly when the other route is dead ($\min=0$) or the better one is certain ($\max=1$). Otherwise the
gain is strict. The gain over the better route $q\ge p$ is $p(1-q)$.

Justin Hart sealed \texttt{P6.SerialParallel} by decision D-026. This session (P6-T004) proves it.

\section{Background and related work}
Series and parallel systems are the two basic structures of reliability theory. A series system works
only if all of its components work. A parallel system works if at least one does. With independent
components their reliabilities are the product formula and the complementary product
formula~\citep{barlowproschan1975}. Both formulas are elementary.

The serial bound has appeared before in this programme. Part~3 of the sealed statement
\texttt{P3.CascadedRecovery} says that for recovery fractions $r_i\in[0,1]$ the virgin share satisfies
$\prod_i(1-r_i)\le 1-r_j$ for every stage $j$. Setting $p_i=1-r_i$ gives exactly part~1 below, over
\texttt{Fin n} rather than an arbitrary \texttt{Finset}. A Lean proof of \texttt{P3.CascadedRecovery}
was given in run \texttt{runs/2026-11-16\_P3} (task P3-T003) of the Viridis Research
repository~\citep{viridisP3T003}. P6 restates the inequality self-contained and does not import it. A
cascade of recovery stages and a chain of corridor links are both products of fractions in $[0,1]$.

We use Lean~4~\citep{lean4} and Mathlib~\citep{mathlib2020}. The graph-theoretic background for the
previous two statements is classical~\citep{diestel2017}.

\section{The model and the statement}
A route is indexed by a finite set \texttt{s : Finset} $\iota$ of links, with link $i$ crossed with probability
$p_i$. Alternative routes are indexed by a finite set \texttt{s} with success probabilities $q_i$. The
sealed definitions are:
\begin{lstlisting}
def serial {ι : Type} (s : Finset ι) (p : ι → ℝ) : ℝ := ∏ i ∈ s, p i

def parallel {ι : Type} (s : Finset ι) (q : ι → ℝ) : ℝ := 1 - ∏ i ∈ s, (1 - q i)

def Statement : Prop :=
  (∀ (ι : Type) (s : Finset ι) (p : ι → ℝ), (∀ i ∈ s, 0 ≤ p i ∧ p i ≤ 1) →
      ∀ j ∈ s, serial s p ≤ p j) ∧
  (∀ (ι : Type) (s : Finset ι) (q : ι → ℝ), (∀ i ∈ s, 0 ≤ q i ∧ q i ≤ 1) →
      ∀ j ∈ s, q j ≤ parallel s q ∧ parallel s q ≤ 1) ∧
  (∀ p q : ℝ, 0 ≤ p → p ≤ 1 → 0 ≤ q → q ≤ 1 →
      (1 - (1 - p) * (1 - q) = max p q ↔ (min p q = 0 ∨ max p q = 1)))
\end{lstlisting}

\begin{theorem}[\texttt{P6.SerialParallel}]\label{thm:sp}
Let $s$ be a finite index set.
\begin{enumerate}
\item If $0\le p_i\le 1$ for all $i\in s$, then $\prod_{i\in s}p_i\le p_j$ for every $j\in s$.
\item If $0\le q_i\le 1$ for all $i\in s$, then $q_j\le 1-\prod_{i\in s}(1-q_i)\le 1$ for every $j\in s$.
\item For $p,q\in[0,1]$: $1-(1-p)(1-q)=\max(p,q)$ if and only if $\min(p,q)=0$ or $\max(p,q)=1$.
\end{enumerate}
\end{theorem}

\paragraph{What the formulas mean, and what they assume.} \texttt{serial} and \texttt{parallel} are just
real-valued formulas, and Theorem~\ref{thm:sp} is a statement about them. The serial product is the
probability that an animal crosses every link \emph{if} the crossings are independent events. The
parallel formula is the probability that at least one route succeeds \emph{if} the routes fail
independently. Independence is a field assumption, not part of the mathematics. It fails, for example,
when two ``alternative'' routes share a link, or when a fire, flood or road closure hits several routes
at once. With positively correlated failures, the true probability that at least one route succeeds is
in general below the independence formula, so the formula is then optimistic. The serial bound in
part~1 needs less. For any events, $\Pr(\bigcap_i A_i)\le\Pr(A_j)$ holds without independence. But the
quantity proved here is the product, and the product equals $\Pr(\bigcap_i A_i)$ only under independence.

\paragraph{Sanity checks in the sealed module.} Two links of $1/2$ in series give $1/4<1/2$. Two
independent routes of $1/2$ in parallel give $3/4>1/2$. At $p=q=1/2$ part~3's equality fails, as neither
side condition holds. A negative control shows that part~1 needs $p_i\le 1$: links $2$ and $1/2$ have
product $1>1/2$.

\section{Proof}
\paragraph{Mathematical proof.} (1) Write $\prod_{i\in s}p_i=p_j\cdot\prod_{i\in s\setminus\{j\}}p_i$.
The second factor lies in $[0,1]$ because each of its factors does, and $p_j\ge0$, so the product is at
most $p_j$. (2) Apply (1) to $1-q_i\in[0,1]$ to get $\prod(1-q_i)\le 1-q_j$, that is,
$q_j\le 1-\prod(1-q_i)$. The product is non-negative, so $1-\prod(1-q_i)\le1$. (3) Without loss of
generality $p\le q$. Then $\max=q$, $\min=p$, and $1-(1-p)(1-q)-q=p(1-q)$. This is zero exactly when
$p=0$ or $q=1$.

\paragraph{Lean proof.} The Lean file follows the same route. The helper \texttt{prod\_le\_factor}
splits off factor $j$ with \texttt{Finset.mul\_prod\_erase}, bounds the rest with
\texttt{Finset.prod\_le\_one} and finishes with \texttt{mul\_le\_mul\_of\_nonneg\_left}. The statement
quantifies over an arbitrary $\iota$\texttt{ : Type} without decidable equality, so the helper opens with
\texttt{classical}. The first compile attempt failed without it. Part~2 applies the helper to
$1-q_i$ and closes with \texttt{linarith}. Part~3 splits on \texttt{le\_total p q}. The forward
direction uses \texttt{mul\_eq\_zero} on $p(1-q)=0$, and the converse substitutes and calls
\texttt{ring}. The complete candidate file follows.

\begin{lstlisting}
import Mathlib
import ViridisStatements.P6.SerialParallel

namespace Candidate

open ViridisStatements.P6.SerialParallel

/-- A product of factors in `[0, 1]` is at most any one of its factors. -/
theorem prod_le_factor {ι : Type} (s : Finset ι) (p : ι → ℝ)
    (hp : ∀ i ∈ s, 0 ≤ p i ∧ p i ≤ 1) (j : ι) (hj : j ∈ s) :
    ∏ i ∈ s, p i ≤ p j := by
  classical
  rw [← Finset.mul_prod_erase s p hj]
  have h0 : 0 ≤ p j := (hp j hj).1
  have h1 : ∏ i ∈ s.erase j, p i ≤ 1 :=
    Finset.prod_le_one (fun i hi => (hp i (Finset.mem_of_mem_erase hi)).1)
      (fun i hi => (hp i (Finset.mem_of_mem_erase hi)).2)
  calc p j * ∏ i ∈ s.erase j, p i ≤ p j * 1 := mul_le_mul_of_nonneg_left h1 h0
    _ = p j := mul_one _

theorem serial_le {ι : Type} (s : Finset ι) (p : ι → ℝ)
    (hp : ∀ i ∈ s, 0 ≤ p i ∧ p i ≤ 1) (j : ι) (hj : j ∈ s) : serial s p ≤ p j :=
  prod_le_factor s p hp j hj

theorem parallel_bounds {ι : Type} (s : Finset ι) (q : ι → ℝ)
    (hq : ∀ i ∈ s, 0 ≤ q i ∧ q i ≤ 1) (j : ι) (hj : j ∈ s) :
    q j ≤ parallel s q ∧ parallel s q ≤ 1 := by
  have hq' : ∀ i ∈ s, 0 ≤ 1 - q i ∧ 1 - q i ≤ 1 := fun i hi => by
    obtain ⟨a, b⟩ := hq i hi; constructor <;> linarith
  have hle := prod_le_factor s (fun i => 1 - q i) hq' j hj
  have hnn : 0 ≤ ∏ i ∈ s, (1 - q i) := Finset.prod_nonneg (fun i hi => (hq' i hi).1)
  unfold parallel
  constructor <;> simp only at hle <;> linarith

theorem two_routes (p q : ℝ) (hp0 : 0 ≤ p) (hp1 : p ≤ 1) (hq0 : 0 ≤ q) (hq1 : q ≤ 1) :
    1 - (1 - p) * (1 - q) = max p q ↔ (min p q = 0 ∨ max p q = 1) := by
  rcases le_total p q with h | h
  · rw [max_eq_right h, min_eq_left h]
    constructor
    · intro e
      have : p * (1 - q) = 0 := by linarith
      rcases mul_eq_zero.mp this with h1 | h1
      · exact Or.inl h1
      · exact Or.inr (by linarith)
    · rintro (e | e) <;> subst e <;> ring
  · rw [max_eq_left h, min_eq_right h]
    constructor
    · intro e
      have : q * (1 - p) = 0 := by linarith
      rcases mul_eq_zero.mp this with h1 | h1
      · exact Or.inl h1
      · exact Or.inr (by linarith)
    · rintro (e | e) <;> subst e <;> ring

theorem main : ViridisStatements.P6.SerialParallel.Statement :=
  ⟨fun ι s p hp j hj => serial_le s p hp j hj,
   fun ι s q hq j hj => parallel_bounds s q hq j hj,
   fun p q hp0 hp1 hq0 hq1 => two_routes p q hp0 hp1 hq0 hq1⟩

end Candidate
\end{lstlisting}

\section{Results}
The candidate \texttt{proofs/P6.SerialParallel.lean} compiles under Lean \texttt{v4.28.0} with Mathlib
commit \texttt{8f9d9cff6bd728b17a24e163c9402775d9e6a365}. It passed the local check script, which runs the
denylist, compile, statement-match, axiom and kernel-replay stages. Its only axioms are \texttt{propext},
\texttt{Classical.choice} and \texttt{Quot.sound}. The independent verifier decides whether the claim is
labelled proven, and its certification page precedes this paper.

\begin{table}[h]
\centering
\begin{tabular}{lll}
\toprule
Configuration & Formula value & Comparison\\
\midrule
Two links $1/2$, $1/2$ in series & $1/4$ & below weakest link $1/2$\\
Links $0.9,0.9,0.9$ in series & $0.729$ & below weakest link $0.9$\\
Two routes $1/2$, $1/2$ in parallel & $3/4$ & gain $1/4$ over best\\
Routes $0.9$ and $0.5$ in parallel & $0.95$ & gain $0.5\cdot0.1=0.05$\\
Routes $0.9$ and $0$ in parallel & $0.9$ & equality (dead alternative)\\
Routes $1$ and $0.5$ in parallel & $1$ & equality (certain route)\\
\bottomrule
\end{tabular}
\caption{Worked examples of Theorem~\ref{thm:sp}, computed by hand from the formulas. The parallel
values are the probability that at least one route succeeds only under independent route failures.}
\label{tab:examples}
\end{table}

\section{Planning reading}
Together with the previous notes, the statement gives a planner three facts about the model. A tree of
corridors is the cheapest connected plan, but every corridor is a bridge. Surviving any single corridor
loss costs at least one more corridor. Along a route, reliability can only fall as links are added, and
by part~1 it is capped by the weakest link. Across routes, under independence, it can only rise, and
part~3 says when a second route adds nothing. The gain $\min(p,q)\,(1-\max(p,q))$ is largest when the
better route is weak and the alternative is strong. A redundant corridor is therefore most valuable
where the existing route is least reliable. This is a property of the formula. Whether an extra corridor
in a real landscape delivers it depends on the independence premise and on actual crossing rates, and
this note claims neither.

\section{Registered test plan: P6-NucleationIslandSize}
The engine asked this session to run the registered plan \texttt{P6-NucleationIslandSize}, which tests
premise \texttt{M-Nucleation} of project P2 on the tree-island recruitment data of Zahawi, Holl, Cole and
Reid~\citep{zahawi2014data}. We copied the frozen code and design file unchanged, and their hashes match
the registration. We downloaded the outcome file from the registered URL (sha256
\texttt{b35a35bb\ldots 861f}) and ran the entry script. It returned \textsc{not\_estimable}, with zero island
sites found. The cause is a column-name mismatch: the frozen code reads the treatment from a column named
\texttt{Treat}, but the file's header calls it \texttt{Trt}. So every row was skipped. We did not change the
analysis. The verdict therefore bears neither for nor against \texttt{M-Nucleation}. A corrected plan would
need a new registration.

\section{Limitations}
\begin{itemize}
\item These are elementary facts of probability and algebra, not new theorems.
\item The proof covers the formulas only. It does not establish that link crossings or route failures are
independent, or any value of a real crossing probability. The premise review filed with this run lists
both independence assumptions.
\item Routes in a real network share links. Then the network's reliability is not a single
series-parallel formula, and computing it exactly is hard in general~\citep{valiant1979}. Nothing here
addresses that case.
\item Part~3 covers two routes. For more routes the analogous equality condition is not sealed.
\item The test-plan run gave no information about its premise, for the reason stated above.
\end{itemize}

\section{Reproducibility}
The candidate proof is \texttt{proofs/P6.SerialParallel.lean}. It imports only \texttt{Mathlib} and the
sealed module \texttt{ViridisStatements.P6.SerialParallel} (sha256 \texttt{a882c978\ldots 89e3}), and it
uses none of the forbidden constructs. To check it, run
\texttt{viridis-lean-check proofs/P6.SerialParallel.lean P6.SerialParallel} under toolchain
\texttt{leanprover/lean4:v4.28.0} and Mathlib commit \texttt{8f9d9cff6bd728b17a24e163c9402775d9e6a365}.
For the plan run, \texttt{code/analyze.py} (frozen) runs from the package root with
\texttt{data/design.json} and \texttt{data/tree\_recruits.csv}, both pinned in \texttt{data\_pins.json}. It
writes \texttt{results/verdict.json}.

\section*{Changelog}
\begin{itemize}
\item 2026-11-23 (P6-T001): first version. Drafted three statement proposals
(\texttt{P6.MinimalNetwork}, \texttt{P6.SingleFailureRobust}, \texttt{P6.SerialParallel}).
\item 2026-11-24 (P6-T002): \texttt{P6.MinimalNetwork} sealed (D-024); mathematical and Lean proofs.
\item 2026-11-25 (P6-T003): \texttt{P6.SingleFailureRobust} sealed (D-025); proof, attainment witnesses,
informal placement observation.
\item 2026-11-26 (P6-T004): \texttt{P6.SerialParallel} sealed (D-026). Paper refocused on route reliability:
statement, independence caveat, link to \texttt{P3.CascadedRecovery}, mathematical and Lean proofs,
worked examples, planning reading. Proof listings of earlier statements removed (they remain in earlier
packages). Report of the \texttt{P6-NucleationIslandSize} plan run added.
\end{itemize}

\bibliographystyle{plainnat}
\bibliography{refs}
\end{document}
